\section{Eventual doubling in the binary-weight array}

For positive integers $i,j$, let $T(i,j)$ be the $j$th positive integer
whose ordinary binary expansion contains exactly $i$ ones, listed in
increasing order. This is the array of OEIS A067576, introduced by
Wilson in 2002. Kimberling conjectured in November 2025 that every
column of $T(i,j)+1$ eventually forms a geometric progression of ratio
two \cite{oeis067576}. We prove the conjecture and determine its exact
starting row.

\begin{theorem}\label{a067576:main}
For positive $i,j$ with $i\ge j-1$,
\begin{equation}\label{a067576:formula}
 T(i,j)+1=(2^j-1)2^{i+1-j}.
\end{equation}
Consequently the least positive row $r$ such that
\[
 T(i+1,j)+1=2\bigl(T(i,j)+1\bigr)\qquad\text{for every }i\ge r
\]
is exactly $\max(1,j-1)$.
\end{theorem}

\subsection*{Exact ranks from binary prefix counts}
Let $w(n)$ be the number of ones in the binary expansion of the natural
number $n$, including $w(0)=0$, and put
\[
 C_i(m)=\#\{n\in\mathbb N: n<m,\ w(n)=i\}.
\]
The identities $w(2n)=w(n)$ and $w(2n+1)=w(n)+1$ imply
\begin{align}
 C_{i+1}(2m)&=C_{i+1}(m)+C_i(m),\label{a067576:even}\\
 C_{i+1}(2m+1)&=C_{i+1}(m+1)+C_i(m).\label{a067576:odd}
\end{align}
Indeed, the even integers below $2m$ correspond to integers below $m$
of unchanged weight; the odd integers correspond to integers below $m$
of weight one less. Below $2m+1$, the even part additionally contains
$2m$, giving the second equation.

Binary induction proves $2^{w(n)}\le n+1$. At an even successor step,
$2^{w(2n)}=2^{w(n)}\le n+1\le2n+1$; at an odd step,
$2^{w(2n+1)}=2\cdot2^{w(n)}\le2n+2$. The base $n=0$ is equality.
Also $w(2^i-1)=i$, since the binary expansion is a string of $i$ ones.
It follows that
\begin{equation}\label{a067576:minimum}
 C_i(m)=0\quad(m<2^i),\qquad C_i(2^i)=1.
\end{equation}
For the first equality, every $n<m$ satisfies $n+1\le m<2^i$ and
therefore cannot have weight $i$. The second follows by adjoining the
single integer $2^i-1$.

Induction in \eqref{a067576:even} now gives
\begin{equation}\label{a067576:nearfull}
 C_i(2^{i+1})=i+1.
\end{equation}
The case $i=0$ counts only zero. For the induction step, the two terms
in \eqref{a067576:even} are $C_{i+1}(2^{i+1})=1$ and
$C_i(2^{i+1})=i+1$. The two largest integers below $2^{i+1}$ have
weights
\[
 w(2^{i+1}-2)=i,\qquad w(2^{i+1}-1)=i+1.
\]
Subtracting their contributions from \eqref{a067576:nearfull} shows
\begin{equation}\label{a067576:base}
 w(2^{i+1}-2)=i,\qquad C_i(2^{i+1}-2)=i.
\end{equation}

\begin{lemma}\label{a067576:rank}
For all $k,d\ge0$, set $v(k,d)=(2^{k+1}-1)2^d-1$. Then
\[
 w(v(k,d))=k+d,\qquad C_{k+d}(v(k,d))=k.
\]
\end{lemma}
\begin{proof}
For $d=0$, this is \eqref{a067576:base}. Suppose it holds for $d$,
and write $v=v(k,d)$. Then
\[
 v(k,d+1)=2v+1,\qquad
 v+1=(2^{k+1}-1)2^d<2^{k+d+1}.
\]
The weight of $2v+1$ is $k+d+1$. Equations
\eqref{a067576:odd} and \eqref{a067576:minimum} give
\[
 C_{k+d+1}(2v+1)
 =C_{k+d+1}(v+1)+C_{k+d}(v)=0+k.
\]
This proves the induction step.
\end{proof}

For every positive $i$, the integers $(2^i-1)2^a$, as $a$ varies, show
that the weight-$i$ row is infinite. Its increasing enumeration is
therefore characterized by membership and the exact number of smaller
members. Lemma~\ref{a067576:rank} says that $v(k,d)$ is the $(k+1)$st
integer of weight $k+d$. For $j\ge1$ and $i\ge j-1$, take
$k=j-1$ and $d=i+1-j$. Then $k+d=i$, and the lemma gives
\eqref{a067576:formula}. Multiplying its right side by two gives the
formula at row $i+1$, proving eventual doubling from
$\max(1,j-1)$ onward.

\subsection*{The starting row is minimal}
For $k\ge0$, consider
\[
 u=2^{k+2}+(2^k-1).
\]
Adding the isolated leading bit shows $w(u)=k+1$. More generally, for
$0\le n<2^r$ one has $w(2^r+n)=1+w(n)$, and hence for
$0\le m\le2^r$,
\begin{equation}\label{a067576:leading}
 C_{i+1}(2^r+m)=C_{i+1}(2^r)+C_i(m).
\end{equation}
The equality splits the count into the integers below $2^r$ and those
having the new leading bit. Applying it with $i=k$, $r=k+2$ and
$m=2^k-1$, and then using \eqref{a067576:nearfull} and
\eqref{a067576:minimum}, yields
\[
 C_{k+1}(u)=C_{k+1}(2^{k+2})+C_k(2^k-1)=k+2.
\]
Thus $u=T(k+1,k+3)$, so
\[
 T(k+1,k+3)+1=5\cdot2^k.
\]
At the next row, the already proved formula gives
\[
 T(k+2,k+3)+1=2^{k+3}-1=8\cdot2^k-1
 \ne10\cdot2^k=2\bigl(T(k+1,k+3)+1\bigr).
\]
For a column $j\ge3$, choosing $k=j-3$ proves failure at row $j-2$.
Any earlier proposed onset would include this transition and therefore
cannot work. For $j=1,2$, doubling already holds from the least positive
row, namely one. This completes Theorem~\ref{a067576:main}.

\paragraph{Formalization and attribution.}
The formal weight is \texttt{Nat.bitIndices.length}, and the source
entry is the zero-based selector \texttt{Nat.nth} at index $j-1$.
The proof establishes infinite positive rows, strict increase, exact
membership, and exact rank, as well as the universal formula and the
\texttt{IsLeast} onset statement. All four source modules compile with
Lean 4.19.0 and pinned mathlib, and their principal theorem audits use
only the standard axioms \texttt{propext}, \texttt{Classical.choice},
and \texttt{Quot.sound}.

Binary-weight ranking is classical. A036563 records the diagonal
special case, and A066884 presents this array with reversed
antidiagonals. The paper of Dobri\'c, Skyers and Stanley develops
polynomial-time weight ranking and unranking \cite{dobric2016}.
The minimum-onset formula was also already posed as an admitted Lean
statement in open-math-agents \cite{oma067576}; that file contains no
proof. We claim a complete proof of the recorded eventual-doubling
assertion and its exact onset, with no claim of first formulation or
global historical priority. The package counts as one source-conjecture
resolution.
